\scanpage{206}
\mspage{62}
est $\mathbb{D}$-coasphérique, d'où $b'\Rightarrow a''$. Il est clair que l'implication inverse est valable si $\alpha_i^*$ est ess. surjectif, c'est le cas $2^\circ$. Or $2^\circ)$ est satisfait si $I$ est préordonné, d'où le cas $1^\circ$. Le cas $4^\circ$ a déjà été traité (4.1.). Il faut encore traiter le cas $3^\circ$, qui contient d'ailleurs le cas $2^\circ$ donc le cas $1^\circ$ (\deleted{\unsure{voir}}\inserted{et aussi} le cas $4^\circ$, cf plus bas).

Il faut exprimer l'isomorphie de la flèche $(\unsure{xx|x})$, pour $F_i$ de la forme $\alpha_i^*(F)$. Soit $\xi$ objet de $A=\mathbb{D}(e)$, il faut exprimer l'isomorphie du $\operatorname{Hom}(\xi,-)$ des deux mbs de $(\unsure{\ast|x})$, i.e. de

\[
 \begin{aligned}
 \operatorname{Hom}(\xi_{I_{/i}},F_i)&\longrightarrow
 \operatorname{Hom}\bigl(\underbrace{\xi'_{I'_{/i}}}_{\deleted{\unsure{u_{/i}^*(\xi_{I_{/i}})}}},u_{/i}^*(F_i)\bigr)\\
 &\unsure{\vcenter{\hbox{\tikz{\node[rotate=90,inner sep=0pt] {$\simeq$};}}}|\shortparallel}\quad\mathrm{adj.}\\
 &\operatorname{Hom}\bigl((u_{/i})_!(\xi'_{I'_{/i}}),F_i\bigr),
 \end{aligned} \tag*{(\unsure{xx|\ast\ast})}
\]

flèche qui se déduit de

\[
 (u_{/i})_!(u_{/i})^*\deleted{\unsure{\xi}}\unsure{\bigl(}\xi_{I_{/i}}\bigr)
 \xrightarrow{\mathrm{adj}}\xi_{I_{/i}}
 \tag*{(\unsure{xxx|\ast\ast\ast})}
\]

en prenant les $\operatorname{Hom}$ dans $F_i$. La condition $u_{/i}$ coasphérique \deleted{\unsure{sig}} (i.e. $b'$) signifie que $(xxx)$ est un iso pour $\forall\ \xi$ de $A$. \uline{Suffit}-il pour ceci que \unsure{ses|son} $\operatorname{Hom}$ dans un $F_i$ soit bijectif, quand on restreint $F_i$ dans $\mathbb{D}(I_{/i})$ à être dans l'``image de $\mathbb{D}(I)$ par $\alpha_i^*$ ? Or cette flèche $(xx)$, pour $F=\alpha_i^*F$, s'explicite par la formule d'adjonction relative
% archive sheet number: 206
